\section{Discussão Final}

\subsection{Comparativo Geral das Abordagens}

O trabalho evoluiu em três estágios distintos, cada um respondendo a uma limitação
do estágio anterior.

\textbf{Serial.}
A implementação sequencial estabeleceu a baseline: uma thread lê o CSV, calcula a
distância de cada ponto aos centróides e acumula os resultados. O perfil de execução
revelou que $\approx$87\% do tempo era gasto em parsing CSV
(\texttt{FloatingDecimal.parseDouble}, \texttt{String.split}), e apenas $\approx$13\%
em computação aritmética. Isso significa que o algoritmo em si é rápido; o gargalo
é a entrada de dados.

\textbf{Concorrente v2 (platform, virtual, hybrid).}
A paralelização da fase de atribuição (lotes de 5.000 pontos, pool de $N$ workers)
não trouxe ganho de tempo por iteração: o parsing CSV seguia sequencial na thread
principal, dominando o tempo total. Os microbenchmarks mostraram B2
$\approx$60 s/op para todas as variantes, idêntico ao serial. O macrobenchmark
revelou throughput de 1,6--4,6 exec/min dependendo do número de threads JMeter
e da variante de GC.

A exploração das três variantes foi útil para entender o comportamento de cada
modelo de thread. Platform threads escalonam pelo kernel; virtual threads
multiplexam sobre carriers do \texttt{ForkJoinPool}. A variante hybrid foi a que
mais claramente separou os papéis: o coordenador virtual desmontava durante I/O e
\texttt{Future.get()}, tornando o parsing CSV invisível ao profiler de CPU
--- revelando que o problema não era a coordenação, mas a leitura em si.

\textbf{Concorrente v3 (dados em memória).}
Carregar o dataset inteiro em \texttt{double[][]} antes do loop de iterações
eliminou o parsing CSV das medições. B2 caiu de $\approx$60 s para $\approx$1{,}6 s
($\approx$37$\times$); o throughput saltou de $\approx$1,6 para $\approx$52
exec/min com 1 thread JMeter. O perfil não registrou nenhum evento de leitura de
arquivo durante o clustering. O custo foi o aumento do heap de $\approx$250 MB
para $\approx$2{,}7 GB.

A Tabela~\ref{tab:comparativo-geral} resume os principais números de cada abordagem.

\begin{table}[H]
\centering
\begin{tabular}{lrrrr}
\hline
\textbf{Abordagem} & \textbf{B2 (s/op)} & \textbf{Throughput 1t} & \textbf{Throughput 16t} & \textbf{Heap pico} \\
\hline
Serial      & 60{,}05 &     ---    &     ---    & $\approx$250 MB \\
v2-platform & 60{,}69 & 1{,}65/min & 4{,}49/min & $\approx$250 MB \\
v2-virtual  & 71{,}61 & 1{,}62/min & 4{,}55/min & $\approx$250 MB \\
v2-hybrid   & 60{,}44 & 1{,}60/min & 4{,}53/min & $\approx$250 MB \\
v3          &  1{,}60 & 52{,}5/min & 62{,}3/min & $\approx$2{,}7 GB \\
\hline
\end{tabular}
\caption{Comparativo geral entre abordagens (B2 = custo por iteração, JMH; throughput = exec/min, JMeter ParallelGC).}
\label{tab:comparativo-geral}
\end{table}

\subsection{Pontos Positivos e Negativos de Cada Abordagem}

\subsubsection*{Serial}

\textbf{Positivos.}
Simples de implementar e depurar. Ausência de condições de corrida por definição.
Baixo consumo de memória. Serve como baseline clara para medir ganhos das versões
concorrentes.

\textbf{Negativos.}
Subutiliza processadores multi-core. A thread principal acumula os dois papéis
(I/O e computação) sem possibilidade de sobreposição.

\subsubsection*{v2 --- Concorrência com Lotes e Futures}

\textbf{Positivos.}
Paraleliza a fase computacionalmente intensiva (atribuição de pontos a centróides)
sem exigir sincronização explícita, graças à acumulação local por worker. A
estrutura de \texttt{Future.get()} fornece garantia de happens-before clara. A
variante hybrid demonstrou que virtual threads são adequadas para coordenadores
que alternam entre I/O e espera.

\textbf{Negativos.}
O parsing CSV sequencial na thread principal continua sendo o gargalo, tornando o
ganho da paralelização invisível no tempo de parede. A criação de objetos
\texttt{List.copyOf} e \texttt{PartialResult} por lote gera pressão de GC.
A virtual thread em tarefas CPU-bound de curta duração introduz variância elevada
no JMH (até $\pm$260 s), pois o \texttt{ForkJoinPool} não foi projetado para
tarefas que nunca bloqueiam.

\subsubsection*{v3 --- Dados em Memória com CountDownLatch}

\textbf{Positivos.}
Elimina o I/O das iterações, revelando o verdadeiro custo computacional do algoritmo.
A divisão em $N$ chunks fixos é mais simples que a criação dinâmica de lotes:
sem \texttt{List.copyOf}, sem \texttt{Future} por lote. \texttt{CountDownLatch}
é mais direto que coletar e iterar sobre uma lista de \texttt{Future}s. O
\texttt{KMeansSampler} com \texttt{sharedPoints} estático permite que múltiplos
threads JMeter compartilhem os dados sem alocar múltiplos GBs por thread.

\textbf{Negativos.}
O dataset inteiro precisa caber na heap ($\approx$1{,}75 GB para $\approx$5{,}75
milhões de pontos com 38 features). Isso inviabiliza a abordagem para datasets
maiores sem particionamento externo. O tempo de inicialização (leitura única do
CSV) não é contabilizado nas métricas, o que pode ser enganoso se o processo for
reiniciado frequentemente.

\subsection{Lições Aprendidas}

\textbf{Medir antes de paralelizar.}
Sem o profile da versão serial, seria natural assumir que paralelizar o loop de
atribuição bastaria para acelerar o algoritmo. O JFR mostrou que 87\% do tempo era
parsing CSV --- tornar paralelo os outros 13\% dificilmente produziria ganho
perceptível.

\textbf{Gargalos migram.}
Depois que o I/O foi eliminado na v3, o verdadeiro gargalo computacional tornou-se
visível: \texttt{squaredDistance} e \texttt{closestCentroid} no loop interno. Em
implementações futuras, esse é o ponto a otimizar (SIMD, vetorização, redução de
dimensionalidade).

\textbf{Virtual threads não são bala de prata para CPU-bound.}
Virtual threads brilham quando há bloqueio I/O: o carrier é liberado e reutilizado.
Em tarefas puramente aritméticas, elas não oferecem vantagem sobre platform threads
e podem introduzir variância de escalonamento. A escolha certa depende do perfil
da tarefa.

\textbf{O modelo de acumulação local escala bem.}
O padrão map-reduce local (cada worker acumula em arrays próprios, a thread
coordenadora reduz no final) evitou toda sincronização interna sem locks, atomics
ou \texttt{volatile}. A corretude foi verificável formalmente via JCStress, que
confirmou a ausência de race conditions nas operações reais e a presença delas
quando a proteção é removida (JCS3, JCS4).

\textbf{O custo de abstração importa.}
A troca de \texttt{List<Future<PartialResult>>} por \texttt{CountDownLatch} + array
de resultados eliminou a alocação de centenas de objetos por iteração. Pequenas
decisões de design (tamanho do lote, estrutura de retorno, tipo de sincronizador)
têm impacto mensurável em cargas de trabalho longas.
